\section{Problem Formulation and Workload Characterization}\label{sec:problem}

\subsection{The context-retrieval workload}
Context-aware trajectory prediction augments each AIS \emph{anchor}---a
position fix $a=(\mathrm{lat},\mathrm{lon},t)$---with a spatial context
window before it is fed to the predictor. Formally, given a static set of
$N$ OSM features $F=\{f_1,\dots,f_N\}$ (each a 2-D MBR with a category
tag) and a stream $S$ of timestamped vessel positions, the pipeline must
compute, for every anchor $a$, the triple
\begin{equation}
\mathcal{C}(a)=\bigl(\,\underbrace{R_F(a,r_1)}_{\text{OSM range}},\;
\underbrace{D(a)}_{\text{SDF}},\;
\underbrace{K_S(a,r_2,k)}_{\text{$k$NN}}\,\bigr),
\end{equation}
where $R_F(a,r_1)=\{f\in F: f.\mathrm{MBR}\cap B(a,r_1)\neq\emptyset\}$ is
the set of features whose MBR intersects the $r_1\!=\!5$\,km query box;
$D(a)\in\mathbb{R}^{H\times W\times 2}$ is the pair of
$128\!\times\!128$ signed distance fields over that patch; and
$K_S(a,r_2,k)$ are the $k$ nearest vessels within $r_2\!=\!3$\,km in the
current snapshot. A training corpus evaluates $\mathcal{C}$ over millions
of anchors, so the workload is \emph{ingest-once, query-many}: $F$ and the
historical $S$ are built once and queried hundreds of thousands of times.

\subsection{Profiling: the context dominates}
Table~\ref{tab:profile} replays each stage of the reference pipeline in
isolation on a $1\,000$-anchor subset. Two facts stand out. First, the
\emph{model is not the cost}: the three retrieval stages account for
essentially the entire per-anchor wall-clock. Second, within them the SDF
stage dominates ($176.4$\,ms/anchor, $50\%$ of the total), and $89\%$ of
that is spent in the quadratic distance transform \texttt{\_udist}; the
neighbour scan adds a further $25\%$. Extrapolated to the
$5.48$-million-anchor corpus, context construction costs
$\approx\!17$ CPU-days---before a single training epoch begins.

\begin{table}[t]
\centering
\caption{Per-stage cost of the reference context pipeline on a
$1\,000$-anchor subset (single-thread CPU). The three retrieval stages
are the entire budget; the SDF stage dominates.}\label{tab:profile}
\small
\begin{tabular}{lrrl}
\toprule
Stage & wall-clock & per-anchor & dominant cost \\
\midrule
OSM tile load + parse & $85.2$\,s  & one-time   & cold file I/O \\
SDF compute           & $176.4$\,s & $176.4$\,ms & \texttt{\_udist} ($89\%$) \\
Neighbour scan        & $88.2$\,s  & $88.2$\,ms  & brute-force scan \\
\bottomrule
\end{tabular}
\end{table}

\subsection{Why brute force is the wrong default}
Each stage is a database query forced into a linear pass. The OSM stage
recomputes a tile id per anchor, reloads JSON, and re-parses it---no index
at all---so it pays $O(N)$ per query and cold I/O on every miss. The SDF
stage runs an $O(HW\!\cdot\!M)$ per-pixel nearest-neighbour loop where the
exact field is computable in $O(HW)$. The neighbour stage rescans the full
snapshot per anchor instead of indexing the moving points. The workload
therefore exhibits three classic database fingerprints---static range,
streaming $k$NN, and output-sensitive distance transform---and the central
question of this paper is how to index it \emph{with correctness
guarantees and without perturbing the downstream model}. We answer it with
the M-CTX framework (Section~\ref{sec:overview}).

\subsection{Design goals and the amplification metric}\label{ssec:goals}
The reformulation fixes three goals that an indexed implementation must
meet before a production pipeline will adopt it. \textbf{(G1) Correctness:}
each stage must return exactly the brute-force answer---for the range
stage we demand \emph{provable} recall, not merely high empirical recall.
\textbf{(G2) Amortised efficiency:} because the workload is
ingest-once/query-many, the figure of merit is the sum of build,
per-query, and footprint costs over the corpus, not query latency in
isolation. \textbf{(G3) Drop-in equivalence:} the emitted context must be
numerically identical to the reference, so downstream predictions are
unchanged. We further track, for the range stage, the \emph{candidate
amplification}
\begin{equation}
\alpha(q)=\frac{|\{\text{features scanned by the exact filter}\}|}{|R_F(q)|},
\end{equation}
the average number of exact-MBR tests performed per true hit. For any
recall-complete index $\alpha\!\ge\!1$, and minimising $\alpha$ at recall
$1.0$ is the central efficiency lever a learned range index controls;
BR-LZ minimises it among recall-complete one-curve indices
(Section~\ref{ssec:extent}). These three goals and the amplification
metric organise the design of M-CTX and the evaluation of
Section~\ref{sec:bench}.
